% ==============================================================================
% Doc. 263 — Fractal Holography: contact points between FFGFT and the
% holographic principle, and the open scaling question at the area law
% Johann Pascher, June 2026
% ==============================================================================

\chapter*{Doc.~263 -- Fractal Holography: FFGFT and the Holographic Principle}
\addcontentsline{toc}{chapter}{Doc.~263 -- Fractal Holography}

\section*{Preliminary remark}

This document locates FFGFT relative to the holographic principle. It
strictly separates three levels: \textbf{(1)}~contact points already
documented in the FFGFT corpus; \textbf{(2)}~an open scaling question at the
area law, computed here as a \emph{reduction layer} but \emph{not} claimed as
a core derivation; and \textbf{(3)}~genuine tensions where FFGFT explicitly
does \emph{not} match the standard picture of holography. The separation is
deliberate: otherwise the word \enquote{holographic} loads an AdS/CFT
expectation that FFGFT neither meets nor intends to meet.

\medskip

\textbf{Methodological status.} Section~\ref{sec:263-doc} reports what is
documented (Doc.~201). Section~\ref{sec:263-open} computes a consequence of
\emph{one assumption} and is marked as a reduction layer. The choice of
assumption is a physical decision, not a derivation.

\medskip

\textbf{Symbols used.}

\noindent
\begin{tabular}{@{}ll@{}}
\toprule
Symbol & Meaning \\
\midrule
$\xi$ & basic geometric parameter, $\xi = 4/30000 \approx 1.333\times10^{-4}$ \\
$\tilde T\cdot m=1$ & time--mass duality (fundamental relation); $\tilde T$ intrinsic time scale \\
$D_f$ & fractal dimension, $D_f = 3-\xi$ \\
$E_0$ & characteristic energy, $E_0=\sqrt{m_e m_\mu}=7.398$\,MeV \\
$E_H$ & Hubble energy, $E_H = E_0\,\xi^{41/4} \approx 1.41\times10^{-33}$\,eV \\
$H_0$ & Hubble quantity, $H_0 = E_H/\hbar \approx 66.2$\,km/s/Mpc \\
$R_H$ & maximal scale, $R_H = \hbar c/E_H \approx 1.4\times10^{26}$\,m \\
$L_0$ & sub-Planck granulation, $L_0 = \xi\,\ell_P \approx 2.2\times10^{-39}$\,m \\
$L_\xi$ & characteristic $\xi$ length ($\sim 100\,\mu$m, Doc.~091) \\
$\ell_P$ & Planck length \\
$z,\,z_{\text{obs}}$ & redshift; $z_{\text{obs}}=\lambda_b/\lambda_e-1$ (wavelength ratio) \\
$\rho_{\text{vac}},\rho_{\text{CMB}}$ & vacuum/CMB energy density; $\rho_{\text{CMB}}=\xi\hbar c/L_\xi^4$ \\
$m_T,\lambda$ & mediator mass $m_T=\lambda/\xi$ (Doc.~201); regulates $\rho\le 1/\xi^2$ \\
$c,\hbar,\hbar c$ & SI conversion factors (not ontological primitives) \\
$T_{00}$ & energy density from $\mathcal{L}=\xi(\partial E)^2$; $T_{00}=\xi[(\partial_t E)^2+(\nabla E)^2]$ \\
\bottomrule
\end{tabular}

\medskip
\noindent\emph{Layers:} core derivation $>$ bridge $>$ reduction layer
(plausibility). Standing caution: $H_0$, $z$, $\Omega$ are $\Lambda$CDM
pipeline outputs, not model-neutral.

% ==============================================================================
\section{Documented contact points}
\label{sec:263-doc}

\subsection*{Area law from node counting (Doc.~201, §7.6)}

FFGFT sets the entropy of a black hole as node configuration entropy:
\[
S = \frac{A}{4\,\ell_P^2} \;\longrightarrow\; S = N_{\text{nodes}}\,\ln 2,
\qquad N_{\text{nodes}} \propto \frac{A}{\xi_0^2},
\]
with $\ell_P^2 \sim \xi_0^2$ in the large limit. This reproduces the
Bekenstein-Hawking area law -- the holographic core statement that entropy
scales with \emph{area}, not \emph{volume}. In FFGFT this is not a postulate;
it falls out of counting saturated nodes on the core surface.

\subsection*{Information as an access problem, not an existence problem}

The winding numbers $(n_\theta, n_\varphi, n_3, n_4)$ are topological
invariants; they cannot disappear. An apparent \enquote{information loss} at
the horizon is, in FFGFT, epistemic (a problem of access), not ontological (a
problem of existence). This is a distinct position on the information
paradox, entering the same terrain as holography -- conservation of
information -- but by a different mechanism: topological conservation rather
than boundary encoding.

\subsection*{Mass as scale projection}

FFGFT already uses the term \enquote{holographic projection}: particle masses
arise as a projection between scales (the VEV onto the sub-Planck scale). This
is holography in the sense of a \emph{projection between scales}, not in the
sense of a bulk/boundary duality.

% ==============================================================================
\section{Open scaling question at the area law (reduction layer)}
\label{sec:263-open}

\subsection*{Decompactification forces corrections -- but not their form}

A natural argument runs: since we live in a decompactified regime (the torus
dimensions are hidden at the scale $L_0 \sim \ell_P/\xi_0$), the hidden
compactness \emph{must} show up as a fractal correction in observable
quantities. This is correct -- and it is precisely $\xi_0$: the correction
parameter is documented in the corpus (mass formula $m_i =
r_i\,\xi_0^{p_i}\,v$; geometric redshift as fractal path-lengthening). That
fractal corrections exist \emph{at all} is therefore not speculation.

\medskip

\noindent
This argument does \emph{not}, however, fix \emph{which} geometric quantity
the correction acts on at a black-hole horizon. \enquote{Decompactification
forces corrections} does \emph{not} entail $S \propto R^{2-\xi_0}$. A
correction on the horizon \emph{area} (Assumption~B below) and one on the
enclosed \emph{volume} (Assumption~C) are both fully decompactified and
fractal -- yet they yield \emph{different} area laws (broken vs.\ exact). The
step from \enquote{$\xi_0$ corrections exist} to \enquote{the horizon entropy
deviates by ${\sim}1\,\%$} is therefore an \emph{additional} assumption about
where the correction acts, not a consequence of decompactification itself.

\subsection*{The question}

Doc.~201 uses the \emph{plain} area law: $N_{\text{nodes}} \propto A \propto
R^2$, exponent exactly $2$. A correction by the fractal dimension $D_f = 3 -
\xi_0$ is \emph{not} built in there. It is therefore an open question whether
$D_f$ affects the horizon \emph{area} at all -- and if so, how.

\subsection*{Three assumptions compared}

\begin{itemize}
  \item \textbf{Assumption A (plain area law, Doc.~201):}
        $S \sim R^{2}$. Exponent exactly $2$.
  \item \textbf{Assumption B (fractal area):} if the nodes sit on a
        codimension-1 boundary of a $D_f$-dimensional structure, the
        effective node count scales as $R^{D_f - 1} = R^{2-\xi_0}$.
        Exponent $2 - \xi_0$.
  \item \textbf{Assumption C (fractal volume):} if $D_f$ acts on the enclosed
        \emph{volume} ($V \propto R^{D_f}$) while the area stays integer, the
        area law remains \emph{exact} at exponent $2$.
\end{itemize}

\subsection*{Result of the computation}

\begin{center}
\begin{tabular}{lcc}
\toprule
Assumption & Exponent & Deviation from area law \\
\midrule
A -- plain area law   & $2$ & $0$ \\
B -- fractal area     & $2 - \xi_0 = 1.99986667$ & $\Delta_{\text{exp}} = \xi_0 = 1.33\times10^{-4}$ \\
C -- fractal volume   & $2$ & $0$ (area law stays exact) \\
\bottomrule
\end{tabular}
\end{center}

\medskip

\noindent
Only Assumption~B breaks the Bekenstein-Hawking law -- by \emph{exactly
$\xi_0$ in the exponent}. The relative deviation grows logarithmically with
size:
\[
\frac{S_B}{S_A} - 1 \;\approx\; -\,\xi_0 \,\ln\!\left(\frac{R}{\ell_P}\right).
\]
The tiny $\xi_0$ is lifted by the logarithm across many orders of magnitude;
for astrophysical black holes the effect reaches percent level:

\begin{center}
\begin{tabular}{lcc}
\toprule
Object & $R/\ell_P$ & Deviation (Assumption B) \\
\midrule
Solar-mass BH ($R\sim 3$\,km)         & $1.9\times10^{38}$ & $-1.17\,\%$ \\
Sagittarius~A* ($R\sim 1.2\times10^{10}$\,m) & $7.4\times10^{44}$ & $-1.37\,\%$ \\
M87* ($R\sim 1.9\times10^{13}$\,m)  & $1.2\times10^{48}$ & $-1.47\,\%$ \\
\bottomrule
\end{tabular}
\end{center}

\subsection*{Interpretation}

This is qualitatively different from the $\xi_0$ effect in other contexts
(e.g.\ the dimension-driven proximity $D/(D+1)\to 3/4$, where $\xi_0$ only
appears at $10^{-5}$). Here $\xi_0$ reaches percent level because
$\ln(R/\ell_P)$ bridges some $38$ orders of magnitude. This is precisely
where a \emph{discriminator} would sit between \enquote{true holography}
(area law exact) and \enquote{fractal near-holography} (exponent $2-\xi_0$).

\medskip

\noindent
\textbf{Status (open).} The choice between Assumption~B and~C is a physical
decision about FFGFT, \emph{not} a derivation from the existing corpus.
Assumption~B is plausible (codimension-1 boundary), but it is not derived in
Doc.~201. As long as the question is open, the percent deviation is the
consequence of \emph{one} assumption, not a prediction of FFGFT. Likewise,
$\ell_P^2 \sim \xi_0^2$ was treated here as a pure prefactor (it cancels in
the ratio); any scale-dependent fractality of $\ell_P$ would add a further
term.

% ==============================================================================
\section{Reconstruction: the hologram analogy and its limits}

\subsection*{Documented: a non-invertible projection (Doc.~193)}

Doc.~193 records that evaluating the mode labels $(r_i, p_i)$ to real numbers
is \emph{not invertible}: once $r_i$ and $p_i$ have been evaluated to numbers
in $\mathbb{R}$, they carry \enquote{no memory} of their origin from
different particles. From a real value (a mass, say, or a derived quantity)
the underlying quantum-number structure cannot be uniquely recovered -- the
fibre over each numerical value has collapsed. This is a documented,
non-invertible projection from the discrete geometric mode space onto
$\mathbb{R}$.

\subsection*{Analogy (metaphor, not derivation): the physical hologram}

This asymmetry can be illustrated by a physical hologram -- explicitly as an
\emph{illustration}, not a derivation. An optical hologram encodes a spatial
scene in an interference pattern on a surface. To turn it back into a visible
image, the surface alone does not suffice: it requires a reference (readout)
beam carrying the \emph{prior information} (wavelength, phase, recording
geometry). And even at the best recording and reconstruction quality, the
visible hologram remains a \emph{faint image} of the original reality -- part
of the information is in principle unrecoverable.

\medskip

\noindent
On this reading, the observable $3{+}1$-dimensional world relates to the
underlying $T^4$ geometry as a reconstructed hologram relates to the recorded
scene: reconstruction requires prior information, and it is lossy in
principle. The observable world is an \emph{incomplete} image of the substrate
-- not because the recording is poor, but because the projection itself
discards information.

\subsection*{Distinction: this is not the non-closure theorem}

This non-invertible projection must be carefully separated from the
\emph{non-closure theorem} (Doc.~193). The non-closure theorem is a
\emph{different}, purely algebraic statement: the set of admissible pairs
$(r_i, p_i)$ is not closed under multiplication, i.e.\ the product of two
distinct modes $(r_i\cdot r_j,\; p_i+p_j)$ corresponds to no particle of the
model. It concerns the \emph{combinability} of modes with one another, not the
\emph{invertibility} of a projection. Both statements come from Doc.~193 but
are distinct and should not be equated.

% ==============================================================================
\section{Reversibility: an assumption of standard holography vs.\ a derived time asymmetry}

\subsection*{Reversibility is an axiom in standard holography}

Standard holography (AdS/CFT) is formulated as a \emph{duality} -- an exact,
invertible dictionary between bulk and boundary. This invertibility rests on
two premises: the \textbf{unitarity} of the underlying quantum theory
(reversible time evolution, information conservation) and the existence of a
\textbf{conformal boundary} on which the exact encoding lives. Both are
posited, not derived. Reversibility is therefore not a measured finding but a
\emph{structural assumption} of the formalism.

\subsection*{FFGFT derives the time asymmetry (Doc.~157)}

FFGFT does not treat the direction of time as an axiom but derives it from
deeper geometry. Doc.~157 anchors the breaking of time-reversal symmetry in
three ways:

\begin{itemize}
  \item \textbf{Algebraic:} $T \cdot m = 1$ with positive-definite mass
        ($m > 0$, antimatter included) forces $T > 0$; a reversal
        $T \to -T$ would require $m \to -m$ and is physically impossible.
  \item \textbf{Topological:} the winding of the vacuum field on the torus has
        a chirality; once chosen, the time direction is topologically
        protected (like a knot that cannot be reversed by continuous
        deformation).
  \item \textbf{Fractal:} at $D_f = 3 - \xi$ the path integrals for
        forward- and backward-running paths differ -- the fractal structure
        turns the discrete symmetry $t \to -t$ into a continuous asymmetry.
\end{itemize}

\noindent
The two situations are therefore \emph{not} symmetric: standard holography
\emph{posits} reversibility, whereas FFGFT \emph{derives} non-reversibility.
Non-invertibility is thus not a separate extra assumption in FFGFT but a
consequence of structures the framework already carries ($T\cdot m=1$, torus
winding, $D_f$).

\subsection*{Empirical limit: only practical experience}

Honesty is required here. The only empirical anchor for the directedness of
time at all is the \textbf{practical experience} of directed happening -- not
a direct measurement and not a theorem. The observed CP/T violation is, as
Doc.~157 itself notes, too weak for the macroscopic arrow of time. FFGFT thus
does not \emph{measure} the asymmetry; it replaces the thermodynamic stopgap
(an unexplained low-entropy initial condition) with a \emph{geometric
derivation} of the directedness already experienced in practice. This is a
structural advantage -- the \enquote{Past Hypothesis} becomes unnecessary --
but it is not an empirical proof of the premise.

\subsection*{Consequence: reversibility as a possible artefact}

From this follows the genuine open question lying beneath both frameworks.
\emph{If} the world is fundamentally time-asymmetric -- as practical
experience suggests and FFGFT derives geometrically -- \emph{then} the exact
reversibility of the holographic mathematics is an \textbf{artefact of the
unitary idealization}: it abstracts away a symmetry breaking that FFGFT
grounds. This is not a claim that standard holography is \enquote{wrong}, but
the observation that its reversibility is an unmeasured assumption -- and that
FFGFT posits a derived asymmetry at the same point. The two assumptions have
different observable consequences; a concrete discriminator has already been
named (the logarithmic entropy deviation $-\xi\,\ln(R/\ell_P)$ under
Assumption~B, Section~\ref{sec:263-open}).

% ==============================================================================
\section{Anchoring at the bounded core (no singularity)}

\subsection*{The anchor}

A projection is admissible as a pure model; but for it to be
\emph{FFGFT-intrinsic} (rather than borrowed from AdS/CFT), it must be
\emph{anchored} at a physical point, like the time-mass duality. T0 has no
singularity: the vacuum density is bounded by
$\rho \le 1/\xi^2 \approx 5.6\times10^{7}$ (Doc.~201), regulated by the
mediator mass $m_T = \lambda/\xi$, which prevents unbounded growth of
$\Delta m$. Where classical GR would run into a singularity, the
$\xi$-bounded core holds against it -- exactly as $T\cdot m=1$ caps the
divergence ($m\to\infty$ is capped by $T\to 0$). The holographic projection is
anchored at this core.

\subsection*{Computational consequence: a cutoff instead of a singularity}

This entails a \emph{physical cutoff}. There is a finite core radius
$R_\text{core}$ below which the nodes are saturated; the area-law scaling
$S \sim R^{a}$ holds only for $R \ge R_\text{core}$, and below it the entropy
saturates. Order of magnitude:
\[
  \frac{R_\text{core}}{\ell_P} \sim \frac{1}{\xi} \approx 7.5\times10^{3}
  \qquad (\log_{10} \approx 3.9).
\]
The $O(1)$ prefactor is model-dependent; only the \emph{order of magnitude}
follows from $\rho \le 1/\xi^2$. Consequence for the computation in
Section~\ref{sec:263-open}: values with $\log_{10}(R/\ell_P) \lesssim 3.9$ lie
\emph{below} the core -- in the region of the non-existent singularity -- and
are physically meaningless. Only above $R_\text{core}$ is the scaling valid.
For astrophysical black holes ($R/\ell_P \sim 10^{38}\text{--}10^{48} \gg
R_\text{core}$) the deviation $-\xi\,\ln(R/\ell_P)$ remains valid unchanged --
the cutoff affects only the unphysical small-scale region.

\subsection*{Interpretive limit: conserved information is geometric base information}

What is conserved at the core -- the node configuration entropy
$N_\text{nodes}\ln 2$ -- is \emph{not} a maximal description of the past
reality, but only the \textbf{base information about the geometric structure}
(winding numbers as topological invariants). Conservation here means
\enquote{access to the geometry}, not \enquote{reconstruction of the
history}. This is consistent with the access-not-existence character
(Section~\ref{sec:263-open}) and with the lossy, directed projection: the
history of the infalling material is not conserved countably, only its
geometric signature.

\subsection*{Open: no exact translation quantity}

A gap remains to be stated. So far we have only a \emph{scale anchoring}
($R_\text{core} \sim \ell_P/\xi$), not a derived \textbf{translation
quantity} that would gauge the area law (core entropy $\leftrightarrow$
horizon area) as an exact map with a fixed coefficient -- analogous, say, to
the central charge in AdS/CFT or to the SI conversion factors $\hbar, c$
within FFGFT itself. No such quantity is documented in the corpus, and none
has been found. This is consistent with the position taken here: no exact
duality, hence no exact translation coefficient -- only a lossy projection
with a scale anchoring. The exact translation quantity remains an open point.



\subsection*{Boundaryless topology}

Standard holography (AdS/CFT) requires a conformal \emph{boundary}. The
universal 4-torus $T^4$ is boundaryless (compact without boundary). A naive
bulk/boundary duality therefore does not apply. FFGFT's
\enquote{holography} is an area law plus scale projection -- not boundary
encoding.

\subsection*{No holographic dark energy}

Models that obtain $\Lambda$ from a holographic horizon are incompatible with
the static FFGFT universe (no $\Lambda$, no dark energy). Here FFGFT stands in
\emph{opposition} to holography, not in contact. An equation of state
$w_{DE}$ does not exist in FFGFT; the comparable observable is the geometric
redshift law (fractal path-lengthening, cf.\ Doc.~041/182), not a
$w$ parameter.

% ==============================================================================
\section{\texorpdfstring{$H_0$ as a projection artefact, not a translation quantity}{H0 as a projection artefact, not a translation quantity}}

\subsection*{No direct \texorpdfstring{$\xi$}{xi}--\texorpdfstring{$H_0$}{H0} bridge -- but a chain through \texorpdfstring{$\hbar c$}{hbar c}}

The opposition to holographic dark energy tempts one to ask whether FFGFT's
$\xi$ and the $\Lambda$CDM parameter $H_0$ can be \emph{translated} into one
another. A \emph{direct} identification $\xi/L_0 = H_0/c$ fails -- but not
because energy and length were incompatible. In natural units ($\hbar = c =
1$) length and energy are dual; the passage to SI metres is exactly one
conversion factor ($\hbar c = 1.973\times10^{-7}$\,eV\,m), and in FFGFT
$\hbar, c$ are SI conversion factors anyway, not ontological primitives. The
error of the naive bridge is not the energy--length duality but the
\emph{choice of scale}: it couples $\xi$ to the sub-Planck granulation
$L_0 = \xi\,\ell_P \approx 2.2\times10^{-39}$\,m instead of to the cosmological
scale $R_H \approx 1.4\times10^{26}$\,m. Both lengths convert correctly into
energies via $\hbar c$ -- they are simply \emph{different} scales (ratio
$R_H/L_0 \approx 6.5\times10^{64}$, which is calibrated to the measured $H_0$, not following from $\xi$ alone).

\subsection*{The correct translation chain: $\hbar c$ as the SI factor}

Through $\hbar c$ the chain does close (Doc.~182):
\[
  \xi \;\longrightarrow\; E_H = E_0\,\xi^{41/4}
  \;\xrightarrow{\;\times\,\hbar c\;}\; R_H = \frac{\hbar c}{E_H}
  \;\longrightarrow\; H_0 = \frac{c}{R_H} = \frac{E_H}{\hbar}.
\]
With $E_0 = \sqrt{m_e m_\mu} = 7.398$\,MeV this gives $E_H = 1.412\times
10^{-33}$\,eV, $R_H = 1.398\times10^{26}$\,m and $H_0^{T0} \approx
66.2$\,km/s/Mpc ($-1.9\%$ from the Planck value). Here $\hbar c$ is precisely
the SI conversion factor that turns the energy representation ($E_H$) into the
length representation ($R_H$) -- exactly the quantity one expects when length
is readable as energy in natural units. So $\xi$ fixes the energy \emph{and} --
via $\hbar c$ -- the associated length; only the \emph{cosmological} length
$R_H$, not the sub-Planck length $L_0$.

\subsection*{The mechanism: $H_0$ is a misread slope}

This chain is at the same time the \emph{explanation} of $H_0$ as an artefact
-- consistent with the projection reading of holography. The mechanism: for
small $z$ the fractal path-lengthening is linear, $z(d) = \kappa\,d$ with
$\kappa = E_H/(\hbar c)$. An observer in the expansion picture reads the same
linear $z(d)$ curve and \emph{names} its slope $H_0/c$. Hence
\[
  H_0^{\text{(measured)}} = \kappa\,c = E_H/\hbar,
\]
i.e.\ $H_0$ is not a dynamical expansion parameter but the \emph{geometric
coefficient} of the path-lengthening, misinterpreted as expansion. The
artefact is the \emph{reinterpretation of a slope}, not a measurement error.
The translation chain $\xi \to E_H \to R_H \to H_0$ does exist -- it only
shows that the translated quantity is geometric, not dynamical. Beyond the
maximal scale $R_H$ one has $z\to\infty$; no $\Lambda$CDM-type particle
horizon is needed.

\subsection*{$H_0$ is emergent, not fundamental}

What matters is where $H_0$ sits ontologically: it is \emph{not a fundamental
constant} of the theory. In the compactified $T^4$ form there is no $H_0$ --
there are no cosmological distances there, no vacuum traversed as a medium, and
$\hbar, \varepsilon_0$ become \enquote{real} only upon decompactification
(Doc.~262). Since $H_0 = E_H/\hbar$ contains this only-then-real $\hbar$, $H_0$
is an \emph{emergent} quantity of the decompactified description -- just as
age, critical density and Hubble length are \emph{Friedmann concepts without a
direct T0 counterpart} (Doc.~016). Only $\xi$ is fundamental; $H_0$ arises at
the transition compact $\to$ decompactified. The non-integer exponent $41/4$
expresses precisely this scale transition (scale flow over $\sim 40$ orders of
magnitude), not a pure geometric constant; forcing an integer reading would
wrongly push an emergent transition quantity back into the compact geometry.
(Earlier formulations calling $H_0$ a \enquote{fundamental constant} or the
redshift an \enquote{energy loss in the vacuum} stem from an outdated stage;
cf.\ Doc.~190, R16.)

\subsection*{The same holds for $z$ itself: real view, misreading, translation}

The reinterpretation concerns not only the slope $H_0$ but the whole $z$
value -- and here the clean three-way split is worthwhile.

\emph{Model-neutral (real):} what is measured is a wavelength ratio,
$1+z_{\text{obs}} = \lambda_{\text{obs}}/\lambda_{\text{emit}}$. This quantity
is observable and identical in both pictures; it is not an artefact.

\emph{The real FFGFT view:} $z_{\text{obs}}$ is the accumulated \emph{fractal
path-lengthening} of the photon over the physical distance $d$,
$z_{\text{obs}} = (L_{\text{eff}}-d)/d$ (Doc.~026). No scale factor, no
expansion, no time evolution of space -- a purely geometric path-length
relation in the static $T^4$.

\emph{The misreading ($\Lambda$CDM):} the same $z_{\text{obs}}$ is read as an
\emph{expansion} redshift, i.e.\ as a cosmological scale factor
$1+z = a_0/a_{\text{emit}}$. In this reading the initial linear slope becomes
$H_0$ and a large value such as $z\approx 1100$ becomes the
\enquote{decoupling} epoch of a recombination. Both are the expansion reading
of the \emph{same} geometric path-lengthening -- not a second, independent
measurement (cf.\ P15/R15, Doc.~190; the static theory has no cosmological $z$
in the expansion sense).

\emph{Translatable:} the two views convert into one another because they
describe the \emph{same} $z_{\text{obs}}(d)$ curve. $H_0$ is the slope at small
$z$, $z\approx 1100$ an extreme value near the maximal scale $R_H$; both are
$\Lambda$CDM readings of the one FFGFT geometry (the FFGFT-own value of this
extremum is $z_*\approx 875$, Doc.~267; $z\approx 1089/1100$ is only the expansion
reading, not the FFGFT value). The translation is therefore
a \emph{reinterpretation of the curve}, not a physical identity of parameters:
what is real is $z_{\text{obs}}$, what is an artefact is solely its expansion
reading.

\subsection*{Placement and open point}

This is precisely the resolution of the apparent dark-energy opposition: a
translation \emph{does} exist -- the chain $\xi \to E_H \to R_H = \hbar c/E_H
\to H_0$ -- but no \emph{direct} $\xi \leftrightarrow H_0$ identity at the
sub-Planck scale (which would obscure its pipeline character and contradict
the standing caution). The model-neutral shared quantity remains the
observable $z(d)$: same curve, two mechanisms, distinguishable at large $z$.
\textbf{Open -- and to be placed:} the exponent $41/4$ is not freely chosen but
measures the \emph{scale transition} from the characteristic energy $E_0$ to
the Hubble energy $E_H$ via the scale-flow/recursion structure of the $\xi$
field (Doc.~026/182, recursion~R Doc.~203). Two things must be separated:
(i)~the \emph{structure} of the chain $\xi\to E_H\to R_H\to H_0$ is documented
and consistent -- $H_0$ is therein a geometric coefficient misread as
expansion, not a dynamical parameter; $\hbar$ acts only as SI conversion
($H_0=E_H/\hbar$), not as a physical input (cf.\ section \enquote{Constants as
projection factors}). (ii)~the \emph{numerical value} of the exponent $41/4$
itself is not yet fully derived from the $\xi$-field theory; Doc.~165 discusses
a unit-free integer form $n=10$ (with several possible geometric readings,
itself open). As long as (ii) is open, the $H_0$ \emph{value} is a reduction
layer -- a strong consistency argument ($-1.9\%$ from the Planck value), not a
closed core result. The \emph{interpretation} ($H_0$ as a projection artefact)
does not depend on (ii) and stands independently.


% ==============================================================================
\section{Constants as projection factors -- and the vacuum energy}

The $H_0$ reading is a special case of a more general principle that recurs
across several quantities and supports the holographic projection idea
quantitatively. It runs in four steps; all four are reproducible as short
calculations and are flagged below as reduction layer.

\subsection*{1. The \enquote{constants} ground nothing -- they convert}

In natural units ($\hbar=c=1$, pure ratios) the FFGFT description is static
and correction-free (Layer~1, Doc.~261). The quantities $c$, $\hbar c$ and
$\sqrt{\hbar G/c^5}$ are not physical inputs but \emph{conversion factors} of
the SI projection (Layer~2):
\begin{itemize}
  \item $c$ is the metre/second factor that arises when the quantities length
        and time -- identical in Layer~1 ($\tilde T\cdot m=1$) -- are projected
        into separate units (Doc.~134: $c=L/T$, \enquote{not a law of nature
        but a ratio}). The value $299\,792\,458$ is the metre definition, not a
        $\xi$ prediction.
  \item $\sqrt{\hbar G/c^5}$ grounds no time scale; it expresses the
        already-fixed scale in seconds (pure unit calibration).
  \item Even $G$ follows $c$-free from $\xi$: $G_{\text{nat}}=\xi^2/(4m_{\text{char}})$
        (Doc.~012/013); $c,\hbar$ enter only at the SI conversion.
\end{itemize}

\subsection*{2. The vacuum structure is fully present when compactified}

The mode structure ($\xi/L_\xi^4$, fractal mode counting at $d=3+\delta$,
Doc.~091) is \emph{fully present} in Layer~1 -- merely unit-free. \emph{Nothing
is added} at decompactification: the same structure appears in Layer~2,
multiplied by the conversion factor $\hbar c$, as a dimensional energy density
$\rho_{\text{vac}}=\xi\,\hbar c/L_\xi^4$ (form from Doc.~091/025). The vacuum
energy thus does not \emph{come into being} -- it merely becomes \emph{visible}
from the SI view through the decompactification of the previously enrolled
fractal/temporal structure.

\subsection*{3. No static energy floor: the energy is dynamical}

That there is no \emph{self-standing} (static) vacuum energy follows
computationally from the energy-field Lagrangian $\mathcal{L}=\xi\,(\partial_\mu
E)(\partial^\mu E)$ with field equation $\Box E=0$ (massless, no potential term
$V(E)$, Doc.~010). The canonical energy-momentum tensor gives
\[
  T_{00} = \xi\left[(\partial_t E)^2 + (\nabla E)^2\right].
\]
$T_{00}$ contains \emph{no} static term. Switching off the temporal component
($\partial_t E=0$) and requiring a source-free vacuum solution
($\Box E=0\Rightarrow\nabla^2 E=0\Rightarrow E=\text{const}$) gives
$T_{00}=0$. The vacuum energy is therefore entirely bound to the temporal
(and, under boundary conditions such as the Casimir effect, spatial)
\emph{variation} of the field -- not an eigenvalue of a field at rest but an
expression of the dynamics. The \emph{free} vacuum has no spatial part; the
\emph{confined} one (Casimir) does, consistent with the free-vs-confined
distinction in Doc.~091.

\subsection*{4. Delimitation: what of this is robust}

Robust are the quantities in which $\xi$ stands \emph{alone} (no fit factors,
no circular exponents): $T_{\text{CMB}}=\tfrac{16}{9}\xi^2 E_\xi$ and
$|\rho_{\text{Casimir}}|/\rho_{\text{CMB}}=\pi^2/(240\xi)\approx308$.

\emph{Clarification on the nature of this $T_{\text{CMB}}$ relation.} It is an
\emph{algebraic} $\xi$ relation, \emph{not a fractal correction}: the factor
$16/9$ is simply $(4/3)^2$, the square of the numerical prefactor of
$\xi = \tfrac{4}{3}\times10^{-4}$ (cf.\ Doc.~025/028/006); $D_f = 3-\xi$ and a
mode count with fractal dimension do \emph{not} enter this derivation. What is
robust is the \emph{ratio} $T_{\text{CMB}}/E_\xi = (16/9)\xi^2$ (Layer~1,
unit-free). The absolute Kelvin value $2.725$\,K follows only via the SI
conversion (embedding price, Doc.~261); the residual deviation of
$\approx 0.9\%$ (Doc.~013) is \emph{not} derived from $\xi$. A correction of the
form $1-C\xi$ that would close these $0.9\%$ exists in the corpus only as a
\emph{numerical approximation} (e.g.\ $C=275/4$), not as a derivation from the
recursion~$\mathcal{R}$ (Doc.~203) -- it remains open.

The ratio given in Doc.~028,
$\rho_{\text{DE}}/\rho_{\text{DM}}=\xi^{\ln(2.5)/\ln\xi}$, requires separate
placement. It yields exactly $2.5$ for \emph{any} $\xi$ and is therefore
circular \emph{as a prediction}. Yet its target value $2.5$
($\Omega_\Lambda/\Omega_{\text{DM}}$) is itself not a model-neutral measurement
but a $\Lambda$CDM-pipeline output. Read \emph{not} as a prediction but as a
\emph{conversion factor} between the FFGFT geometry and the $\Lambda$CDM
$\Omega$ bookkeeping -- analogous to $c$ (metre/second) and to $z$ -- the
circular calibration via $2.5$ is admissible. But only \emph{partially}: unlike
$c$, where both sides share the same real counterpart, the $\Lambda$CDM side
carries with $\rho_{\text{DE}}$ a model artefact (no $\Lambda$ in FFGFT); the
factor translates partly something real (DM bookkeeping $\leftrightarrow$ $\xi$
effect), partly an artefact. Cf.\ Doc.~190, R17.

The \emph{interpretation} is thereby clear: dark energy is eliminated (no
$\Lambda$, static); dark matter is \emph{not a substance} but a $\xi$-geometric
effect -- an artefact of the $\Lambda$CDM view, like $H_0$ and $z$.
\emph{Quantitatively} open is only the model-neutral reproduction of what
$\Lambda$CDM books as DM ($\Omega_{\text{DM}}$, rotation curves).

\paragraph{Status.} This section is reduction layer: it makes a principle
documented in Doc.~261/262 (Layer~1 correction-free / Layer~2 SI projection)
explicit for $c$, $\rho_{\text{vac}}$ and the dynamics. It predicts no new
cosmological value; the numerical-value circularity of $L_\xi$ (fixed from
$\rho_{\text{CMB}}$) remains a limit.

% ==============================================================================
\section{A formal bridge to UIFT -- verifiable on one side only}

The same projection picture yields a \emph{formal} bridge to the information
quantity $\xi_{\text{UIFT}}$ of Onur Teker's UIFT (Doc.~257). It is placed here,
not asserted: its force is \emph{asymmetric}.

The bridge rests solely on the \emph{first order} $T_{\text{CMB}}/E_\xi =
(16/9)\xi$ (not on the fractal correction $K_{\text{frak}}$ and not on the
absolute Kelvin value). Inserting it into the UIFT relation
$\xi_{\text{UIFT}} = k_B T_{\text{CMB}}\ln 2/(m_e c^2)$ gives a \emph{fixed}
factor
\[
  \frac{\xi_{\text{FFGFT}}}{\xi_{\text{UIFT}}}
  = \frac{m_e c^2/\text{eV}}{(16/9)\ln 2} \approx 414\,684 .
\]
This value is not an empirical result but a pure \emph{conversion factor}
between a Layer-1 quantity ($\xi_{\text{FFGFT}}$, dimensionless) and a
Layer-2 quantity tied to the CMB temperature ($\xi_{\text{UIFT}}$); it
coincides with its closed form.

\textbf{What this verifies -- and what it does not.} The \emph{FFGFT side} is
fully reproducible: $\xi$ is geometric, $(16/9)\xi$ follows from the mode count,
$m_e$ from the lepton hierarchy -- anyone can recompute the factor. The
\emph{UIFT side} we \emph{cannot} verify: the relation
$\xi_{\text{UIFT}}=k_B T_{\text{CMB}}\ln 2/(m_e c^2)$ comes from Teker's work,
and whether it is correct or empirically validated lies beyond our reach. The
bridge is therefore \emph{formal}: \emph{if} Teker's relation holds \emph{and}
the FFGFT values hold, \emph{then} a fixed factor links the two quantities.
What is thereby verified is the FFGFT side of the bridge -- not the UIFT side.

% ==============================================================================
\section{Summary}

\begin{itemize}
  \item \textbf{Documented:} FFGFT reproduces the Bekenstein-Hawking area law
        from node counting (Doc.~201); information is topologically conserved
        (access, not existence problem); mass as scale projection.
  \item \textbf{Open (reduction layer):} if $D_f = 3-\xi_0$ affects the
        horizon area (Assumption~B), the exponent becomes $2-\xi_0$ with a
        logarithmically growing deviation $-\xi_0\ln(R/\ell_P)$ (percent level
        for real BHs). If $D_f$ acts on the volume (Assumption~C), the area
        law stays exact. The choice is to be made on physical grounds.
  \item \textbf{Reconstruction:} the projection of mode labels onto
        $\mathbb{R}$ is non-invertible (Doc.~193) -- illustrated by the
        physical hologram (metaphor): reconstruction needs prior information
        and remains a faint image. Strictly separated from the non-closure
        theorem.
  \item \textbf{Anchoring:} the projection is anchored at the
        $\xi$-bounded core ($\rho \le 1/\xi^2$, no singularity; Doc.~201) --
        like $T\cdot m=1$. Consequence: a cutoff at $R_\text{core}\sim\ell_P/\xi$;
        below it the entropy saturates, the $R^{2-\xi}$ scaling holds only
        above. Real black holes lie far above (deviation valid). Conserved
        information = geometric base information, not the history. Open: no
        exact translation quantity, only a scale anchoring.
  \item \textbf{Reversibility:} standard holography \emph{posits}
        invertibility (unitarity + boundary); FFGFT \emph{derives} the time
        asymmetry (Doc.~157: $T\cdot m=1$, torus winding, $D_f$). The only
        empirical anchor of directedness is practical experience; under that
        premise the exact reversibility is an artefact of the unitary
        idealization.
  \item \textbf{Tension:} $T^4$ is boundaryless (no AdS/CFT duality); no
        $\Lambda$, no $w_{DE}$ (opposition to holographic dark energy).
  \item \textbf{$H_0$ artefact:} no \emph{direct} $\xi\leftrightarrow H_0$
        equation at the sub-Planck scale -- the naive form fails on the choice
        of scale ($L_0=\xi\ell_P$ instead of $R_H$), not on the energy--length
        duality (in natural units both are dual via $\hbar c$). But the chain
        $\xi\to E_H=E_0\xi^{41/4}\to R_H=\hbar c/E_H\to H_0=E_H/\hbar$ does exist
        ($\approx 66.2$ km/s/Mpc; Doc.~026/182), with $\hbar c$ as the SI factor.
        $H_0$ is thus a geometric coefficient misread as expansion; the
        model-neutral quantity remains the observable $z(d)$.
        Open: exponent $41/4$ not yet fully derived.
  \item \textbf{Constants as projection:} $c$, $\hbar c$, $\sqrt{\hbar G/c^5}$
        ground nothing; they convert to SI (Doc.~134/261). The vacuum structure
        is fully present when compactified (unit-free) and becomes SI-visible
        only on decompactification. $T_{00}=\xi[(\partial_t E)^2+(\nabla E)^2]$
        (from $\mathcal{L}=\xi(\partial E)^2$, $\Box E=0$, Doc.~010) has NO static
        term -- without field variation $T_{00}=0$: no self-standing energy,
        the energy comes from the temporal component.
\end{itemize}

\noindent
\textbf{Terminological caution.} \enquote{Holographic} in FFGFT is to be read
as \emph{area law + scale projection}, not as \emph{boundary encoding}. This
distinction should remain explicit in all FFGFT texts to avoid false AdS/CFT
expectations.
